\documentclass[11pt]{article}

% ------------------------------------------------
% Basic packages
% ------------------------------------------------
\usepackage[margin=1in]{geometry}
\usepackage{amsmath, amssymb}
\usepackage{graphicx}
\usepackage{booktabs}
\usepackage{url}
\usepackage{hyperref}
\usepackage{xcolor}
\usepackage{setspace}
\usepackage{natbib}

% Optional: nicer fonts
% \usepackage{lmodern}

% ------------------------------------------------
% Title info
% ------------------------------------------------
\title{Routing Strategies for Multi-Agent LLM Systems:\\
Semantic, Performance-Weighted, and LLM Meta-Routing}

\author{
  Zhaoheng Li \\
  Department of Electrical and Computer Engineering \\
  University of Toronto \\
  \texttt{zhaoheng.li@mail.utoronto.ca}
}

\date{December 2025}

% ------------------------------------------------
% Document
% ------------------------------------------------
\begin{document}

\maketitle

\begin{abstract}
Multi-agent systems powered by Large Language Models (LLMs) offer significant advantages through specialized domain expertise, but they introduce a critical challenge: efficiently routing queries to the most appropriate agent. This paper investigates three complementary routing strategies for LLM-based multi-agent systems: (1) semantic similarity using pre-trained embeddings, (2) performance-weighted adaptive routing that learns from historical feedback, and (3) LLM-based meta-classification that leverages language model reasoning for routing decisions. We evaluate these methods on a labeled dataset of 200 queries across four specialist domains: Statistics, Machine Learning, Data Engineering, and Visualization. Our results demonstrate that adaptive performance-weighted routing achieves 73\% accuracy compared to 69\% for semantic similarity, while maintaining computational efficiency. Furthermore, Method 2's ability to learn during both training and testing phases enables continuous improvement over time. These findings provide practical guidance for practitioners designing multi-agent systems and highlight the importance of combining semantic understanding with empirical performance feedback.
\end{abstract}

% ------------------------------------------------
\section{Introduction}

The rapid expansion of Large Language Models (LLMs) has enabled new paradigms for intelligent systems. Rather than relying on a single monolithic model, modern AI systems increasingly employ multiple specialized agents, each fine-tuned or prompted for expertise in a particular domain. This multi-agent approach offers several advantages: domain experts can be optimized independently, system capacity can scale with the number of agents, and failures in one domain do not necessarily degrade overall system performance.

However, multi-agent systems introduce a fundamental challenge: \emph{query routing}. When a user submits a query, the system must decide which agent to consult. This is non-trivial because queries often span domain boundaries, agents have overlapping expertise, and the consequences of misrouting (poor answer quality, wasted computation) are significant. A routing mechanism must balance several competing objectives: (1) \emph{accuracy}, routing to agents that will provide correct, high-quality answers; (2) \emph{efficiency}, minimizing computational overhead for routing decisions; and (3) \emph{adaptability}, improving routing quality over time as the system gains experience.

The objective of this work is to systematically investigate how different routing strategies balance these objectives. We implement and evaluate three distinct approaches: a semantic similarity baseline, which is computationally lightweight but does not leverage historical performance; a performance-weighted adaptive method that learns per-agent accuracy and self-assessment metrics; and an LLM-based meta-classification approach that uses language model reasoning for routing decisions. By comparing these methods on a common evaluation dataset, we aim to provide empirical evidence for practitioners designing multi-agent systems.

Our contributions are: (1) a concrete implementation and evaluation of three complementary routing strategies; (2) a methodology for adaptive learning in routing that combines external accuracy with agent self-assessment (IN\_DOMAIN classification); (3) empirical results showing that adaptive routing outperforms static baselines while remaining computationally practical; and (4) analysis of the accuracy-latency tradeoff across routing methods.

% ------------------------------------------------
\section{Related Work}

\subsection{Multi-Agent Coordination and Routing}

The coordination of multiple AI agents has become increasingly important as Large Language Models (LLMs) continue to improve. The challenge of routing queries to appropriate specialists has attracted significant research attention. Foundational work on federated learning \cite{Brendan2016} established principles for distributed learning without centralizing data, providing a foundation for multi-agent architectures. More recently, Panda et al.\ \cite{Panda2025} propose adaptive LLM routing under budget constraints, demonstrating that routing decisions can be optimized simultaneously for accuracy and computational efficiency. Their work shows that performance-weighted routing, where agents are selected based on historical track records, can significantly reduce wasted computation while maintaining or improving answer quality.

A closely related direction is confidence-aware routing. Barrak et al.\ \cite{Barrak2025} introduce CARGO, a framework that explicitly incorporates agent self-assessment into routing decisions. If an agent reports confidence that it can handle a query (analogous to our IN\_DOMAIN signal), it becomes more likely to be selected. Zhang et al.\ \cite{Zhang2025} propose Meta-Router, which bridges two evaluation paradigms: gold-standard evaluations based on external ground truth and preference-based evaluations that reflect agent self-assessment. Their key insight—that both signals are valuable—directly motivates our dual-feedback approach combining accuracy and IN\_DOMAIN rates.

\subsection{Semantic Embeddings and Embedding-Based Selection}

Embedding-based methods provide a practical, training-free foundation for routing. Sentence transformers have proven effective for semantic similarity tasks, enabling queries and agent descriptions to be represented in a shared embedding space. Li and Zhou \cite{Li2024} demonstrate that mixture-of-experts LLMs contain implicit embedding representations, suggesting that semantic structure captured by embeddings is meaningful for expert selection. Sridharan et al.\ \cite{Sridharan2025} propose GradientSpace, an unsupervised clustering method that shows semantic embeddings can effectively organize and classify data without labeled training data. These results motivate our Method 1 (semantic baseline), which offers a strong, training-free baseline that serves as a comparison point for more sophisticated methods.

\subsection{Learnable Routing and Adaptive Systems}

A fundamental insight from recent work is that routing quality can improve through learning from experience. El Hili et al.\ \cite{ElHili2025} treat LLMs as in-context meta-learners, showing that models can dynamically select appropriate tools and strategies based on task characteristics. This principle directly underlies Method 2: the router learns which agents are most reliable for different query types and adapts weights accordingly. Similarly, Zhuang et al.\ \cite{Zhuang2025} investigate learnable dynamic routing for mixture-of-experts systems (LD-MoLE), demonstrating that routing weights can be continuously refined based on task performance.

The insight of combining multiple feedback signals—both objective performance metrics and subjective self-assessments—has roots in multi-view learning. Han et al.\ \cite{Han2021} introduce trusted multi-view classification, which dynamically integrates evidence from multiple sources to improve both accuracy and robustness of predictions. We apply this principle to routing by treating external accuracy (ground truth labels) and internal agent assessment (IN\_DOMAIN classification) as complementary signals that together inform routing decisions.

\subsection{Benchmarking and Evaluation Standards}

Establishing rigorous evaluation protocols is essential for progress in multi-agent routing. Wang et al.\ \cite{Wang2025} introduce LoRALib, a standardized benchmark for evaluating expert routing methods in mixture-of-experts systems. Their work establishes best practices for measuring both routing accuracy and computational efficiency. Critically, they emphasize the importance of evaluating on diverse query distributions, distinguishing between straightforward queries (where most routing methods succeed) and challenging, ambiguous queries where intelligent routing provides the greatest benefit. We adopt this principle by evaluating our methods on queries varying in difficulty and clarity.

This paper brings together three complementary routing strategies—semantic similarity, performance-weighted learning, and LLM meta-reasoning—providing a systematic empirical comparison. Our work extends prior research by explicitly combining external ground-truth feedback with agent self-assessment, demonstrating a balanced approach to learning in multi-agent systems.

% ------------------------------------------------
\section{Methodology}

\subsection{System Architecture}

The routing system consists of four specialist agents and a central router that selects the most appropriate agent for each query. All agents are backed by the same base large language model, served locally via Ollama (e.g., \texttt{llama3.1}), but are differentiated by prompts and output schemas.

\singlespacing

Each agent targets a distinct but overlapping sub-domain:

\begin{itemize}
    \item \textbf{StatisticsAgent} handles statistical inference, hypothesis testing, confidence intervals, regression analysis, distribution tests, power and sample size calculations, and experimental design.
    \item \textbf{MLAgent} focuses on machine learning model training and evaluation, including deep learning, feature engineering, cross-validation, hyperparameter tuning, overfitting diagnosis, and model selection.
    \item \textbf{DataEngAgent} covers data engineering tasks such as SQL and schema design, ETL pipelines, data cleaning, warehousing, performance tuning, and data integration.
    \item \textbf{VizAgent} addresses data visualization, including chart selection, dashboard design, exploratory data analysis, and communication of quantitative findings.
\end{itemize}

All agents return a uniform, structured response with specified fields: \texttt{domain} (IN\_DOMAIN or OUT\_OF\_DOMAIN), \texttt{task\_type}, \texttt{difficulty}, \texttt{answer}, and optional metadata (files analyzed, code executions, and visualizations) and latency. This unified schema allows the router and evaluation scripts to treat all agents uniformly.

The router receives a query with optional file paths, and the router then selects one primary agent. It supports three different routing strategies:

\begin{enumerate}
    \item Method 1: semantic similarity baseline.
    \item Method 2: performance-weighted adaptive routing.
    \item Method 3: LLM meta-classification.
\end{enumerate}

The same four agents and the same evaluation dataset are used for all methods.

\subsection{Routing Methods}

\subsubsection{Method 1: Semantic Similarity Baseline}

The semantic similarity method uses sentence embeddings to route queries. A pre-trained sentence transformer (such as \texttt{all-MiniLM-L6-v2}) generates 384-dimensional embeddings.

For each agent $a$, a short natural-language capability description is encoded into an embedding vector $\vec{a}$. At runtime, a query $q$ is encoded as $\vec{q}$. The router computes cosine similarity
\begin{equation}
    \text{semantic}(q,a) =
    \frac{\vec{q} \cdot \vec{a}}{\lVert \vec{q} \rVert \,\lVert \vec{a} \rVert},
\end{equation}
and selects the agent with the highest similarity. Agent embeddings are computed once at initialization; routing cost per query is dominated by a single embedding computation and a few dot products. Method~1 does not rely on historical performance data and requires no pre-training.

\subsubsection{Method 2: Performance-Weighted Adaptive Routing}

Method~2 extends the semantic baseline by incorporating learned per-agent performance statistics. The goal is to bias routing decisions toward agents that have historically been accurate, self-consistent, and fast.

\singlespacing
For each agent $a$, the router maintains:

\begin{itemize}
    \item an \emph{accuracy rate} $\text{accuracy}_a$: the fraction of routed queries where $a$ matched the ground-truth agent label;
    \item an \emph{IN\_DOMAIN rate} $\text{indomain}_a$: the fraction of routed queries where $a$ reported the query as IN\_DOMAIN;
    \item an \emph{average latency} $\ell_a$: an exponential moving average of response time in seconds.
\end{itemize}

Latency is converted into a bounded score
\begin{equation}
    \text{latency\_score}_a = \frac{1}{1 + \ell_a},
\end{equation}
so that faster agents receive higher scores.

For each query $q$, semantic similarity $\text{semantic}(q,a)$ is computed as in Method~1. The final routing score for agent $a$ is
\begin{equation}
    \begin{split}
    S_a(q) ={} & 0.20 \cdot \text{semantic}(q,a)
       + 0.35 \cdot \text{accuracy}_a \\
       & + 0.35 \cdot \text{indomain}_a
       + 0.10 \cdot \text{latency\_score}_a.
    \end{split}
\end{equation}

The weights were chosen empirically to balance query-specific semantics (20\%) with global reliability, giving equal importance to external correctness (35\% accuracy) and the agent's own domain judgement (35\% IN\_DOMAIN), while modestly rewarding lower latency (10\%).

\subsubsection*{Why IN\_DOMAIN Matters}

Including the IN\_DOMAIN classification as a routing signal addresses a fundamental insight: agents have better visibility into their own capabilities than external evaluators. A well-calibrated agent's self-assessment is immediate and cost-free, making it a valuable signal complementary to ground-truth evaluation. While external accuracy captures routing correctness, IN\_DOMAIN rates capture agent reliability—whether the agent's self-knowledge is trustworthy. By giving both signals equal weight, we recognize that these decorrelated feedback sources together improve routing decisions more effectively than either alone.

After each query routed using Method~2, performance statistics for the selected agent are updated via an exponential moving average with learning rate $\alpha = 0.1$:
\begin{equation}
    \text{rate}^{\text{new}} =
    (1 - \alpha)\,\text{rate}^{\text{old}} + \alpha\,\text{feedback},
\end{equation}
where \texttt{feedback} is either $1$ or $0$ depending on whether the event (correct routing, IN\_DOMAIN report, or latency observation) is considered successful. This allows the system to adapt to performance trends while maintaining stability.

\subsubsection{Method 3: LLM Meta-Classification}

The LLM meta-classification method elevates routing itself to an LLM decision. Rather than relying only on fixed embeddings and scalar statistics, a separate call is made to an LLM that sees both the agent descriptions and the user query.

The meta-classifier constructs a system prompt listing each agent and its capabilities. An instruction is set to return only the name of the most suitable agent, and the user query is provided as a separate message. The LLM (e.g., \texttt{llama3.1}) returns free-text output, which is then normalized and mapped to one of the four valid agent names. This method can exploit rich language understanding and task reasoning, but it introduces substantial overhead because each query requires an extra LLM processing call before the primary agent is invoked.

\subsection{Training and Evaluation Protocol}

\subsubsection{Datasets}

Two labeled datasets are used:

\begin{itemize}
    \item A \emph{training set} of 120 queries for pre-training Method~2's performance statistics. These queries are distributed across the four agents (approximately 30 per agent) and cover a range of typical tasks for each domain.
    \item A \emph{test set} of 200 queries for final evaluation. Each test query has an associated ground-truth agent label. The set includes both clear single-domain prompts and boundary cases that could reasonably fit multiple agents.
\end{itemize}

\subsubsection{Pre-Training for Method 2}

Before testing Method~2, the performance-weighted classifier is pre-trained on the 120 training queries. During this phase, each training query is routed to its ground-truth agent, so the accuracy rate approaches $1.0$. The IN\_DOMAIN rate and average latency reflect the agent's actual behavior on these representative queries. These statistics provide initial priors for the test phase.

Importantly, Method~2 continues to update its performance statistics during the test phase, incorporating feedback from the agent's IN\_DOMAIN/OUT\_OF\_DOMAIN classification and the routing accuracy. This enables continuous learning: the router's weights adapt as it processes the 200 test queries, refining its routing decisions based on observed agent performance.

\subsubsection{Testing Procedure and Metrics}

For each routing method, the same 200 test queries are processed. For each query, the router:

\begin{enumerate}
    \item selects an agent according to the chosen method;
    \item invokes the corresponding agent and collects its structured response;
    \item records whether the routed agent matches the ground-truth label (routing accuracy), whether the agent reported IN\_DOMAIN, routing time, agent execution time, and total latency.
\end{enumerate}

Method~2 continues updating its performance statistics during testing, allowing it to learn from agent responses and improve its routing decisions over time.

\subsection{Experimental Setup}

All experiments are executed on a local desktop machine with a 13th~Gen Intel Core i5--13600K processor, 32~GB of RAM, and an NVIDIA GeForce RTX~4070~SUPER GPU. The software stack uses Python~3.11 with \texttt{sentence-transformers} for embeddings, \texttt{pydantic} for schema validation, and Ollama for serving the LLM model (\texttt{llama3.1}). A random seed is set to ensure deterministic query ordering and reproducible runs.

% ------------------------------------------------
\section{Results}

The evaluation comprised five complete test runs with different random seeds (22, 32, 52, 62, 82) to ensure reproducibility and stability. Each run evaluated all three routing methods on 200 test queries across four specialist domains, with consistent ground-truth labels across seeds. A spot-check quality evaluation on 20 sampled responses validated the correlation between routing accuracy and answer quality.

\subsection{Routing Accuracy and Stability}

Method~1 (semantic similarity) consistently achieved 69.0\% accuracy (±0\%) across all trials, demonstrating perfect stability. Method~2 (performance-weighted) ranged from 71.0\% to 75.0\%, with mean 72.8\% (±1.5\%), representing a 3.8 percentage-point improvement over baseline. Method~3 (LLM meta-classification) achieved 94.0\%--95.5\%, with mean 94.7\% (±0.6\%), demonstrating superior accuracy and remarkable consistency. Relative to baseline, Method~2 gained an average of 8.6 correct routes per trial, while Method~3 gained 51.8 routes, highlighting the substantial benefit of LLM reasoning despite computational cost.

In-domain rates reveal increasing confidence with method sophistication: Method~1 reported 91.8\% (±0.5\%), Method~2 reported 94.2\% (±0.8\%), and Method~3 reported 99.2\% (±0.2\%). This pattern indicates semantic routing makes conservative assessments, adaptive learning moderately increases acceptance rates, and LLM routing exhibits near-perfect domain recognition. Method~3's near-perfect in-domain rate shows excellent alignment between self-assessment and actual correctness.

\subsection{Latency and Computational Efficiency}

Routing overhead varied dramatically across methods. Method~1 incurred 11.5ms per query (embedding + cosine similarity), Method~2 added 12.0ms (weighted scoring), while Method~3 required 2410.6ms per routing decision—approximately 200× the overhead of semantic routing due to an additional LLM call before specialist agent invocation.

End-to-end latency per query averaged: Method~1 = 5.49s, Method~2 = 6.80s (24\% increase), Method~3 = 8.24s (50\% increase). The accuracy-efficiency tradeoff is favorable at high-accuracy regimes: Method~2 achieves 3.8pp accuracy gain at 24\% latency cost (0.16pp per percent latency), while Method~3 achieves 25.7pp gain at 50\% latency cost (0.51pp per percent latency), demonstrating superior efficiency for quality-critical applications.

\subsection{Quality Evaluation}

Spot-check evaluation on 20 sampled responses validated that routing accuracy correlates with answer quality. Mean quality score was 4.47/5.0 (σ=0.42), with 90\% rated Good or Excellent. Distribution: 55\% Excellent (4.5--5.0), 35\% Good (4.0--4.5), 10\% Acceptable (3.0--4.0), 0\% Poor. No responses fell below 3.0/5.0, indicating agents maintain professional standards even for misrouted queries.

By agent: StatisticsAgent led at 4.71/5.0 (n=7), followed by VizAgent (4.67/5.0, n=3), with DataEngAgent and MLAgent at 4.25/5.0 and 4.23/5.0 (n=4, 6). The modest 0.48-point spread confirms consistent quality across agents. By dimension: Relevance (4.60/5.0) and Domain Confidence (4.50/5.0) were strongest, while Correctness and Actionability (both 4.40/5.0) showed room for improvement. The highest-quality response (5.0/5.0) was a StatisticsAgent explanation of ANOVA with post-hoc testing; the lowest (3.60/5.0) was an MLAgent response lacking business-specific use cases.

Figure~\ref{fig:quality_breakdown} illustrates quality variation by agent and evaluation dimension. StatisticsAgent and VizAgent consistently outperform baseline, while relevance and domain confidence emerge as agent strengths, suggesting future improvements should deepen correctness and actionability.

\subsection{Summary Results}

Table~\ref{table:complete_results} consolidates metrics across all five trials. Method~1 provides a stable, low-latency baseline (69.0% accuracy, 5.49s per query). Method~2 offers moderate improvements (72.8% accuracy, 6.80s per query) with continuous learning capability, suitable for balanced accuracy-efficiency requirements. Method~3 achieves highest accuracy (94.7%) at highest latency cost (8.24s per query), appropriate for quality-critical scenarios where computational budget allows.

All methods demonstrated high reliability across seeds without catastrophic failures. Method~1's zero variance and Method~3's 0.6pp standard deviation validate robustness for deployment across varied data distributions.

\begin{figure}[h]
\centering
\includegraphics[width=0.95\textwidth]{quality_breakdown.png}
\caption{Quality evaluation breakdown. (Left) Average quality scores by agent, showing StatisticsAgent (4.71/5.0) and VizAgent (4.67/5.0) outperforming baseline. (Right) Quality by dimension: Relevance and Domain Confidence (4.60, 4.50) are strengths, while Correctness and Actionability (4.40) are improvement areas.}
\label{fig:quality_breakdown}
\end{figure}

\begin{table}[h]
\centering
\small
\begin{tabular}{lcccccc}
\toprule
\textbf{Method} & \textbf{Acc.} & \textbf{In-Dom.} & \textbf{Route (ms)} & \textbf{Total} & \textbf{Qual.} \\
\midrule
Semantic & 69.0±0.0 & 91.8±0.5 & 11.5 & 5.49s & 4.41 \\
Performance & 72.8±1.5 & 94.2±0.8 & 12.0 & 6.80s & 4.50 \\
LLM & 94.7±0.6 & 99.2±0.2 & 2410.6 & 8.24s & 4.53 \\
\bottomrule
\end{tabular}
\caption{Five-trial summary: mean and standard deviation. Routing is per-query overhead; Total is end-to-end latency; Quality is spot-check (n=20).}
\label{table:complete_results}
\end{table}

\section{Discussion}

The results demonstrate that routing strategy choice involves a clear accuracy-latency tradeoff with important practical implications. Method~1's semantic similarity baseline, achieving 69.0\% accuracy with minimal overhead, provides a strong foundation for routing tasks. Its training-free nature and perfect stability across seeds (0\% variance) make it attractive for rapid deployment and resource-constrained environments. However, the 31\% misrouting rate indicates substantial room for improvement, particularly on boundary-case queries where semantic similarity alone is insufficient.

Method~2's performance-weighted adaptive routing improves accuracy to 72.8\% (±1.5\%) with modest computational cost (1.31 seconds additional latency across 200 queries). The continuous learning capability—updating statistics both during training and testing—represents a key advantage: the router improves dynamically as it gains experience with agent behavior. The 1.5pp standard deviation across five seeds suggests reasonable stability, though slightly higher variance than Method~1 reflects the adaptive nature of learning-based approaches. Importantly, Method~2's 94.2\% in-domain rate (versus 91.8\% for Method~1) indicates that incorporating agent self-assessment as a routing signal meaningfully improves domain recognition without requiring ground-truth labels in deployment.

Method~3's LLM meta-classification achieves near-perfect accuracy (94.7\% ± 0.6\%), representing a 25.7 percentage-point improvement over baseline. This comes at significant computational cost: 2.4 seconds of routing overhead per query translates to 2.75 hours per 1000-query workload—overhead that is prohibitive for interactive applications. However, the exceptional accuracy and 99.2\% in-domain rate validate that LLM reasoning can capture subtle domain boundaries and contextual cues that simpler methods miss. For quality-critical domains (clinical decision support, financial analysis, legal research), this accuracy premium may justify the latency investment.

Relative to related work, our findings align with Panda et al.'s observations about budget-constrained routing: Method~2 achieves a favorable efficiency frontier, gaining 3.8pp accuracy at only 24\% latency cost (0.16pp per percent latency). This compares favorably to naive approaches that might increase computational budget 10-fold for marginal gains. Our explicit combination of external accuracy and internal domain confidence (IN\_DOMAIN) extends prior work by Zhang et al. on meta-evaluation, providing empirical validation that decorrelated feedback sources improve routing more effectively than either alone.

The quality evaluation (Figure~\ref{fig:quality_breakdown}) provides additional validation: routing accuracy positively correlates with answer quality (mean 4.47/5.0, 90\% good/excellent). Interestingly, agent quality variance is modest (0.48-point spread), suggesting that even misrouted queries receive acceptable answers—agents maintain professional standards across domains. This finding has practical implications: while routing accuracy matters, system robustness does not require perfect routing, creating space for lighter-weight routing methods in non-critical applications.

Limitations of this work merit acknowledgment. First, evaluation is restricted to four agents and 200 test queries; scaling to larger agent ensembles and diverse query distributions remains unexplored. Second, the in-domain ground truth is binary (matching the pre-specified correct agent); boundary-case queries with multiple valid agents may reveal limitations of our evaluation protocol. Third, all agents share a single LLM backbone; heterogeneous agents (different model families, external APIs) may exhibit different routing dynamics.

\section{Conclusion}

This paper evaluates three routing strategies for multi-agent LLM systems. The semantic similarity baseline (Method~1) provides a training-free starting point (69\% accuracy, 5.49s latency). Performance-weighted adaptive routing (Method~2) improves to 72.8\% accuracy with continuous learning and modest computational cost. LLM meta-classification (Method~3) achieves 94.7\% accuracy but at significant latency cost (8.24s), suitable only for quality-critical applications.

A key contribution is the dual-signal routing methodology that combines external accuracy with agent self-assessment (IN\_DOMAIN rates), demonstrating that decorrelated feedback sources improve routing more effectively than either alone.

Future work should explore hybrid routing strategies, scaling to larger agent ensembles, online learning extensions, and heterogeneous agent architectures. These findings demonstrate that multi-agent systems can effectively balance specialized domain expertise with computational efficiency through intelligent routing design.



% ------------------------------------------------
% References
% ------------------------------------------------
\bibliographystyle{plainnat}
\begin{thebibliography}{99}

\bibitem{Barrak2025}
Barrak, A., Fourati, Y., Olchawa, M., Ksontini, E., \& Zoghlami, K. (2025).
\newblock CARGO: A framework for confidence-aware routing of large language models.
\newblock In \textit{Proceedings of the 35th IEEE International Conference on Collaborative Advances in Software and Computing} (pp. 1--12).
\newblock \url{https://arxiv.org/abs/2509.14899}

\bibitem{Brendan2016}
Brendan, M., McMahan, B., \& Ramage, D. (2016).
\newblock Communication-efficient learning of deep networks from decentralized data.
\newblock In \textit{Proceedings of the 20th International Conference on Artificial Intelligence and Statistics} (pp. 1273--1282). PMLR.
\newblock \url{https://arxiv.org/abs/1602.05629}

\bibitem{ElHili2025}
El Hili, Y. A., Thomas, A., Tiomoko, M., Benechehab, A., L\'eger, C., Ancourt, C., \& K\'egl, B. (2025).
\newblock LLMs as in-context meta-learners for model and hyperparameter selection.
\newblock \textit{arXiv preprint arXiv:2510.26510}.
\newblock \url{https://arxiv.org/abs/2510.26510}

\bibitem{Han2021}
Han, Z., Zhang, Z., \& Wang, X. (2021).
\newblock Trusted multi-view classification with dynamic evidence integration.
\newblock \textit{IEEE Transactions on Pattern Analysis and Machine Intelligence}, 43(5), 1501--1515.
\newblock \url{https://doi.org/10.1109/TPAMI.2020.3045810}

\bibitem{Li2024}
Li, Z., \& Zhou, T. (2024).
\newblock Your mixture-of-experts LLM is secretly an embedding model for free.
\newblock \textit{arXiv preprint arXiv:2410.10814}.
\newblock \url{https://arxiv.org/abs/2410.10814}

\bibitem{Panda2025}
Panda, P., Magazine, R., Devaguptapu, C., Takemori, S., \& Sharma, V. (2025).
\newblock Adaptive LLM routing under budget constraints.
\newblock In \textit{Findings of the Association for Computational Linguistics: EMNLP 2025} (pp. 1--15).
\newblock \url{https://arxiv.org/abs/2508.21141}

\bibitem{Sridharan2025}
Sridharan, S., Ravikumar, D., Raghunathan, A., \& Roy, K. (2025).
\newblock GradientSpace: Unsupervised data clustering for improved instruction tuning.
\newblock \textit{arXiv preprint arXiv:2512.06678}.
\newblock \url{https://arxiv.org/abs/2512.06678}

\bibitem{Wang2025}
Wang, S., Lu, Y., Li, Y., Gao, Y., Nie, J., Yu, S., Tian, Y., \& Xuan, Q. (2025).
\newblock LoRALib: A standardized benchmark for evaluating LoRA-MoE methods.
\newblock \textit{arXiv preprint arXiv:2509.18137}.
\newblock \url{https://arxiv.org/abs/2509.18137}

\bibitem{Yurdem2024}
Yurdem, B., \"Ozkaya, E., Ert\"urk, I., \& Kaya, Y. (2024).
\newblock Federated learning: Privacy-preserving distributed machine learning.
\newblock \textit{ACM Computing Surveys}, 56(12), 1--45.
\newblock \url{https://doi.org/10.1145/3626730}

\bibitem{Zhang2025}
Zhang, Y., Xie, F., Yang, S., \& Wu, C. (2025).
\newblock Meta-router: Bridging gold-standard and preference-based evaluations in large language model routing.
\newblock \textit{arXiv preprint arXiv:2509.25535}.
\newblock \url{https://arxiv.org/abs/2509.25535}

\bibitem{Zhuang2025}
Zhuang, Y., Shen, Y., Bian, Y., Su, Q., Ji, S., Shi, Y., \& Miao, F. (2025).
\newblock LD-MoLE: Learnable dynamic routing for mixture of LoRA experts.
\newblock \textit{arXiv preprint arXiv:2509.25684}.
\newblock \url{https://arxiv.org/abs/2509.25684}

\end{thebibliography}

\end{document}
